\subsection{Construction of the actual top lift on every valid cut}
\label{actual-top-lift}

Put $R=r+1$ and $d=2r+k+1$.  Retain the original real-place data,
finite compact representations, actual radial operators and their
normalizations.  Let $1\le l_0\le k$ and
$\sigma=(-1)^{k-l_0}$.  A word is always an original cut word:
its real parameters satisfy the specified cut inequalities, and
distinct parameters have distance greater than two.  The cut
admissibility theorem therefore applies with every required adjacent
shift; no inverse below is taken outside its prescribed domain.
The construction supplies the valid-cut statement in the source,
whose unrestricted differently named version omits the bounds on
$l_0$.  The local contract permits $k=1$, where the compact
Capelli product below is empty; its proof includes that case.
Every arithmetic place package contains these cut bounds and
the additional hypothesis $k\ge2$.\cite{top-lift-contract-source}

Write $D=D_R$, $D_r$ for the indicated leading determinants, and
\[
 \mathcal G=|D|^2-TT'|D_r|^2,\qquad
 \mathbf G^+=\{g:\mathcal G(g)>0\}.
\]
The regular radial coordinates satisfy
$D=D_r\zeta$ and $\mathcal G=|D_r|^2G$, where
$G=\zeta\bar\zeta-TT'$.  These formulas will only identify
operators on the regular cell: the constructed group operators
are smooth on all of $\mathbf G^+$, including $D_r=0$.

\begin{theorem}\label{actual-top-lift-theorem}
For the data above there is a top-lift family $\Theta_{w,e}$ with
the original regular radial action
\begin{equation}\label{actual-top-lift-radial}
 \mathcal J(\Theta_{w,e}F)
   =\bigl[\mathcal F_w\mathcal B_w\bigr]_{e,\mathbf0}
                                                     \mathcal JF.
\end{equation}
Here $\mathcal F_w$ is the prescribed recursive terminal frame,
$\mathcal B_w$ is the original tail-first product of the site
matrices, and $\mathbf0$ is their fixed initial tensor index.
All signs, shifted arguments and rational corrections in those
operators are retained.  The identity applies to every original
smooth compact coefficient; $\mathcal J$ uses the original Haar
measures and leading-minor gauge.
Equivalently the right side may be applied to any smooth covariant
radial function agreeing with $\mathcal JF$ on the regular domain:
all operators there are local, so equality on that open domain
also gives equality of their derivatives.

For $m=|w|$, each component is a finite differential operator in
the original class $\mathrm{InDDPlus}$ with
\begin{equation}\label{actual-top-lift-cost}
                 o=md,\qquad \kappa=11mRd.
\end{equation}
In particular it satisfies the original $\mathrm{InDDplus2}$
condition with the absolute constant $C_b=11$.  Only $D^{-1}$,
$\bar D^{-1}$ and $\mathcal G^{-1}$ occur as scalar denominators;
there is no inverse of an intermediate leading minor.
The empty-word operator is the identity on the entire coefficient
space.  These operators preserve the original smooth rapidly
decreasing class supported on $\mathcal G\ge\varepsilon>0$.
\end{theorem}

\begin{proof}
We give the analytic generators first, then the word construction.
This also makes clear why the inverse finite-type operator does
not introduce an unproved uniform spectral gap.

\paragraph{The actual fibre generators.}
Use the cofactor vector $\Omega(g)$ constructed in
\TopLiftPriorReference{Section}{rl-section}{cofactor construction}.
Its coordinates are polynomials of
degree at most $R$ in the matrix entries and their conjugates,
and $\|\Omega(g)\|^2=\mathcal G(g)/T'$.  Put
$\nu(g)=\Omega(g)/\|\Omega(g)\|$ on $\mathbf G^+$.
The compressed compact Capelli matrix $C_\nu(Y)$ is the actual
Capelli operator on $\nu^\perp$, with parameter
$Y+(k-2)/2$.  The result on compressed Capelli matrices in
\TopLiftPriorReference{Section}{rl-section}{cofactor construction} proves its
phase invariance, compact covariance and adjoint identity.
\TopLiftPriorReference{Proposition}{rl-actual-terminal-lift}{actual terminal lift proposition}
gives precisely
\begin{equation}\label{actual-top-lift-fibre}
 x(Y)\Phi=\Phi\,\sigma C_\nu(Y),\qquad
 x(Y)=\sigma G^{-1}t_{11}(Y),
\end{equation}
where the last matrix uses indices $0,1$.  Thus $x(Y)$ is lifted
by right multiplication by $\sigma C_{\nu(g)}(Y)$.

For the original common finite type set $\mathcal S$, put
$p_\rho(Y)=\prod_{i=1}^{k-1}(Y+\rho_i)$.  At a prescribed
invertible parameter, define
\[
 P_Y(X)=\prod_{\rho\in\mathcal S}(X-\sigma p_\rho(Y)),\qquad
 q_Y(X)=\frac{1-P_Y(X)/P_Y(0)}{X}.
\]
The numerator is divisible by $X$.  The field of right
multiplications by $q_Y(\sigma C_{\nu(g)}(Y))$ therefore lifts
the exact $\xi(Y)=x(Y)^{-1}$ on the original covariant radial
image, by \TopLiftPriorReference{Theorem}{rl-regular-inverse}{regular inverse theorem}.
Repeated spectral
values cause no problem.  No invertibility assertion is made
on unused right constituents.

Here are the bounds needed for these fields, including order
zero.  For either matrix function $H(v)$ just used, dependence
on $v\ne0$ is smooth and homogeneous of degree zero.  Hence
\[
 \|D^aH(v)\|\le C_a\|v\|^{-a},\qquad
 \|L\Omega(g)\|\le C_i(1+\|g\|_{\rm F})^R
                       \quad(|L|\le i).
\]
The second inequality follows by differentiating the fixed
polynomials with left or right invariant unit vector fields;
their degree does not increase.  The ordered chain rule gives
\[
 \|L(H\circ\Omega)(g)\|
   \le C_{\varepsilon,i}(1+\|g\|_{\rm F})^{Ri}
                \quad(|L|=i,\ \mathcal G(g)\ge\varepsilon).
\]
At $i=0$ use the maximum of $H$ on the unit sphere.  Since
$\|g\|_{\rm F}\ge1$ on the unitary group, this implies exactly
the required exponent $3Ri$ for any two unit Lie words whose
total length is $i$.  Smoothness, covariance and the same
adjoint bounds follow from the cofactor construction.  This
is the source's full $\mathrm{FibBound}$, with constants allowed
to depend on the fixed representations, $Y$, $\varepsilon$
and $i$.  In particular no uniform bound as $P_Y(0)\to0$
has been used.\cite{top-lift-fibre-source}

\paragraph{The differential class.}
Say that a radial operator has cost $(o,\kappa)$ if it has a
group lift commuting with the prescribed compact and unipotent
actions, with the stated radial identity, and is a finite sum
of terms
\begin{equation}\label{actual-top-lift-class}
 P(g)D^{-a}\bar D^{-b}\mathcal G^{-c}
       \prod_j(R_{v_j}E_j)(g)\ R_uF(g),
 \quad \deg P\le\kappa,\quad |u|+\sum_j|v_j|\le o.
\end{equation}
The product of fibre endomorphisms has its written order;
each $E_j$ satisfies $\mathrm{FibBound}$.  This is exactly
the original class with an empty set of additional minors.
Addition and scalar multiplication preserve a common cost.
The product rule proves the composition rule
\begin{equation}\label{actual-top-lift-compose}
 (o_1,\kappa_1)(o_2,\kappa_2)
  \longmapsto(o_1+o_2,\kappa_1+\kappa_2+2Ro_1).
\end{equation}
Indeed derivatives of a polynomial do not increase its degree;
each derivative falling on a scalar inverse contributes a
numerator of degree at most $2R$, and a derivative falling on
a fibre factor is recorded in its word $v_j$.  This proves
the assertion for every ordered product without commuting
fibre matrices.  The group covariance is preserved by
composition.  Compact support is preserved by locality, so
the radial identities compose as well.

The preceding fields give cost $(0,0)$ for $x(Y)$ and
$\xi(Y)$.  Every polynomial, finite product or matrix in
these fields also has cost $(0,0)$, irrespective of its
polynomial degree.  Multiplication by $D,\bar D$ has cost
$(0,R)$, and by $D_r,\bar D_r$ cost $(0,r)$.
The multipliers $D^{-1},\bar D^{-1},\mathcal G^{-1}$ have
cost $(0,0)$.  Thus the radial expression
$D_r^{-1}\bar D_r^{-1}G^{-1}$ is lifted as the single
smooth multiplier $\mathcal G^{-1}$; its factors are not
lifted separately.

The signed full Capelli determinant
$\sigma\mathcal D_V(Y)$, with parameter $Y+(d-1)/2$,
has constant invariant-derivative coefficients and order at
most $d$, hence cost $(d,0)$.  Its exact radial image is
\begin{equation}\label{actual-top-lift-central}
 C(Y)=A_r^\vee(Y)n_0(Y)+B_r^\vee(Y)n_1(Y),\qquad
 n_i(Y)=Q_{i0}(Y)A_r(Y)+Q_{i1}(Y)B_r(Y),
\end{equation}
where $Q=\sigma t$ is the actual terminal matrix.  This
follows by applying the bordered Capelli identity of
\TopLiftPriorReference{Lemma}{rr-radial-capelli}{bordered radial Capelli identity}
once for each isotropic pair.
At each step its gauge is the prescribed
$|c_j|^{2(r-j)+k}$, so its shift is exactly the one displayed.
The right and left site products retain their original
order.  Compact quotient properness justifies differentiation
of the Jacquet integral at each step, and the central
Capelli operator commutes with both group actions.
This proves the full group-to-radial identity, not merely a
principal-symbol identity.\cite{top-lift-generators-source}

\paragraph{Four basic operators.}
Put
\[
 b_+=\bar\zeta^{-1}n_1,\quad
 b_-=\zeta^{-1}(Q_{01}A_r^\vee+Q_{11}B_r^\vee),\quad
 a_+=\zeta^{-1}A_r,\quad a_-=\bar\zeta^{-1}A_r^\vee.
\]
All operators in each formula have the common parameter $Y$.
The terminal commutations proved in
\TopLiftPriorReference{Section}{re-section}{recursive frame construction}
and their adjoints give
\[
 [Q_{10},\zeta]=-Q_{11},\quad
 [Q_{01},\bar\zeta]=Q_{11},\quad
 [Q_{00},\zeta]=-Q_{01},\quad
 [Q_{00},\bar\zeta]=Q_{10}.
\]
The crossed commutators for $Q_{10},Q_{01}$ and both
commutators for $Q_{11}$ vanish.  The two $Q_{00}$ identities
use the terminal determinant relation and invertibility at
$Y$ and one adjacent parameter.  All $Q_{ij}$ commute with
the site operators.  Conjugation by $D_r$ shifts every
unbarred Euler derivative by $1$, and conjugation by
$\bar D_r$ shifts every barred Euler derivative by $1$.
Consequently induction on $s\ge0$ gives
\begin{align}
 \bar D^{-s}C(Y)\bar D^s
 &=A_r^\vee(Y-s)\{n_0(Y)+s b_+(Y)\}
                         +B_r^\vee(Y-s)n_1(Y),\label{actual-top-lift-conjplus}\\
 D^{-s}C(Y)D^s
 &=\{A_r^\vee(Y)Q_{00}(Y)+B_r^\vee(Y)Q_{10}(Y)-s b_-(Y)\}
                                                    A_r(Y+s)\nonumber\\
 &\quad+\{A_r^\vee(Y)Q_{01}(Y)+B_r^\vee(Y)Q_{11}(Y)\}
                                                    B_r(Y+s).
                                                    \label{actual-top-lift-conjminus}
\end{align}
For clarity, one conjugation sends $(n_0,n_1,b_+)$ to
$(n_0+b_+,n_1,b_+)$ in the first formula.  In the second
it sends the first brace before its $s$ term to itself
minus $b_-$, and fixes the other brace and $b_-$.  This
proves the induction while preserving the multiplication order.

The site recurrence makes $A_r,A_r^\vee$ monic of degree $r$,
and $B_r,B_r^\vee$ of degree at most $r-1$ (zero when $r=0$).
Both right sides are polynomials in $s$ of degree at most
$R$, with leading coefficients $(-1)^rb_+$ and $-b_-$.
The shifted site polynomials need no additional inverse-domain
assumption at $Y\pm s$.
Set $\lambda_s=(-1)^{R-s}/(s!(R-s)!)$ for $0\le s\le R$.
Lagrange coefficient extraction therefore lifts
\[
 b_+=(-1)^r\sum_{s=0}^R\lambda_s\bar D^{-s}C(Y)\bar D^s,
 \qquad
 b_-=-\sum_{s=0}^R\lambda_s D^{-s}C(Y)D^s.
\]
The cost is $(d,3Rd)$: the numerator degree is at most
$R^2+2Rd\le3Rd$.  The same terminal commutations give
\begin{align}
 a_+(Y+1)&=\bar D\mathcal G^{-1}\xi(Y)
       \{b_+(Y)D_r-D_rD^{-1}b_+(Y)D\},\label{actual-top-lift-ap}\\
 a_-(Y-1)&=-D\mathcal G^{-1}\xi(Y)
       \{b_-(Y)\bar D_r-\bar D_r\bar D^{-1}b_-(Y)\bar D\}.
                                                    \label{actual-top-lift-am}
\end{align}
For example the first brace, after factoring $D_r$, is
$\bar\zeta^{-1}\zeta^{-1}Q_{11}(Y)A_r(Y+1)$;
use $Q_{11}=Gx(Y)$ and $D=D_r\zeta$ to cancel the
displayed factors.  The second identity has the opposite
sign from $[Q_{01},\bar\zeta]=Q_{11}$.
Applying these formulas at the adjacent arguments gives
cost $(d,8Rd)$ for $a_\pm(Y)$, when $Y-1,Y,Y+1$ are
in the original inverse set.  The second terms inside
the braces cost at most $(d,6Rd)$; their left scalar
factor adds $R$, which is within $8Rd$.
All denominators are still the three allowed full-group
ones.\cite{top-lift-basic-source}

\paragraph{The word slopes have bounded shape.}
Let $\mathcal K$ be the algebra generated by $x(Y)$ for
real $Y$ and by $\xi(Y)$ at real invertible parameters.
It is closed under adjoints and all its elements have
cost $(0,0)$.  The recursive formulas in
\TopLiftPriorReference{Theorem}{re-recursive-frame}{recursive frame theorem}
show that every entry of
$x_w(Y),\xi_w(Y),c_w,c_w^{-1}$ belongs to $\mathcal K$ at
the indicated real parameters.  This assertion includes
every old-word rational factor and every shifted scalar
$q_w$.

For a fixed matrix operator $L$, let
\[
 \mathscr S(L;\phi,\tau)=\operatorname{span}_{\mathbb C}
  \{A[L,B]C,\ \phi C,\ \phi\tau C:
                       A,B,C\in\operatorname{Mat}(\mathcal K)\},
 \quad \phi=\frac{\bar\zeta}{\zeta G},\quad
 \tau=\frac{TT'}{\bar\zeta}.
\]
Scalar diagonal embeddings into the word tensor space are
understood.  The two scalars commute with $\mathcal K$, so
this span is stable under left and right multiplication
by matrices over $\mathcal K$.
Write $L(Y)=\ell_\varnothing(Y)$ and
$\delta_w(Y)=\ell_w(Y)-L(Y)$.  We claim
\begin{equation}\label{actual-top-lift-lshape}
 \delta_w(Y)\in\mathscr S(L(Y);\phi,\tau)
           \quad(Y,Y+1\hbox{ in the current inverse set}).
\end{equation}
Here is the full induction, with fresh matrix units
$e_{ij}$.  At a plus step of parameter $a$, set
$X=x_w(a)$ and $X^{-1}=\xi_w(a)$.  The exact RTT relation
identifies the divided difference of the old $e$ entry
with $X^{-1}e_w(Y)X-e_w(Y)$.  Substitution in its two-by-two
update gives
\begin{align}
 \delta_{+}(Y)=\delta_w(Y)&+\phi\tau e_{00}
             -\frac{\phi}{Y-a}e_{01}\nonumber\\
 &+(e_{11}-e_{00})
    \{X^{-1}[L(Y),X]+X^{-1}\delta_w(Y)X-\delta_w(Y)\}.
                                                    \label{actual-top-lift-plusshape}
\end{align}
The connection part commutes with $X$; expanding the
braces directly proves this equality.  At a minus step
of parameter $a$ the same RTT relation cancels the lower
diagonal divided difference, giving exactly
\begin{equation}\label{actual-top-lift-minusshape}
 \delta_-(Y)=\delta_w(Y)+\phi\tau e_{11}
 -\phi\frac{q_w(a-1)}{a-Y-1}\,
                      e_{10}\xi_w(a)\xi_w(a-1).
\end{equation}
In particular the $a-1$ shift, both inverse factors and
the original scalar correction have all survived.
Each displayed term is in the required span, proving
\eqref{actual-top-lift-lshape} from $\delta_\varnothing=0$.

The minus slope is not obtained by an unweighted adjoint.
Its exact definition is
\[
 m_w(Y)=d_w^{-1}\ell_w(Y-1)^*d_w,\qquad
 d_w=\omega^{|w|}c_w,\qquad \omega=\zeta\bar\zeta/G.
\]
Put $M(Y)=m_\varnothing(Y)=L(Y-1)^*$,
$\phi^*=\zeta/(\bar\zeta G)$ and $\tau^*=TT'/\zeta$.
The original connection identities
$[\nabla^*,\bar\zeta]=-1$, $[\nabla^*,\zeta]=0$ and
$[\nabla^*,G]=-\zeta$ give
\[
 \omega^{-1}M(Y)\omega=M(Y)-\phi^*\tau^*.
\]
Indeed $[\nabla^*,\omega]=TT'\zeta/G^2$; the remaining
terms of $M(Y)$ commute with $\omega$.
Taking adjoints maps the span in
\eqref{actual-top-lift-lshape} at $Y-1$ to
$\mathscr S(M(Y);\phi^*,\tau^*)$.
Every member of this latter span commutes with $\omega$:
this follows for the commutators by the Jacobi identity,
since $[M,\omega]=-\omega\phi^*\tau^*$ commutes with
$\mathcal K$.  Thus, for real $Y$ in the stated domains,
\begin{align}
 m_w(Y)-M(Y)
  &=c_w^{-1}[M(Y),c_w]
     +c_w^{-1}\delta_w(Y-1)^*c_w-|w|\phi^*\tau^*\nonumber\\
  &\in\mathscr S(M(Y);\phi^*,\tau^*).
                                                    \label{actual-top-lift-mshape}
\end{align}
This proves the two slope assertions for words of arbitrary
length, with the original weighted adjoint.\cite{top-lift-shape-source}

\paragraph{One letter and every word.}
For a plus letter of parameter $Y$, its exact two entries are
\[
 a_+(Y),\qquad s\{\ell_w(Y)V+W\},\qquad
 s=\bar\zeta^{-1}G,\quad V=A_r(Y),\quad W=B_r(Y).
\]
The initial second entry is
$u=s\{L(Y)V+W\}=\xi(Y)b_+(Y)$, of cost $(d,3Rd)$.
All of $s,V,W,\phi,\tau$ commute with $\mathcal K$.
The three types of slope corrections have the exact reductions
\[
 sA[L,B]CV=A[u,B]C,\qquad s\phi CV=a_+C,\qquad
                      s\phi\tau CV=\tau a_+C.
\]
The first has cost at most $(d,7Rd)$ by
\eqref{actual-top-lift-compose}.  For the other two, the
following scalar reduction avoids an unnecessary derivative
cost.  Since $V,W$
commute with every matrix in $\mathcal K$, so do $a_+$
and $\tau$.  Therefore $a_+C=Ca_+$ and
$\tau a_+C=C\tau a_+$: left multiplication by $C$
has no derivative cost.  The multiplier
$\tau=TT'\bar D_r/\bar D$ has cost $(0,r)$.
Thus these two terms cost $(d,8Rd)$ and $(d,8Rd+r)$,
respectively, within $(d,9Rd)$.
For a minus letter use
$s=\zeta^{-1}G$, $V=A_r^\vee(Y)$, $W=B_r^\vee(Y)$,
$M(Y)$, $\phi^*,\tau^*$ and $u=\xi(Y)b_-(Y)$.
The weighted identity \eqref{actual-top-lift-mshape}
gives the identical estimates, preserving all operator orders.
Both components of every letter therefore have cost $(d,9Rd)$,
independently of the old-word length.

Finally the actual tail-first site product and the recursive
frame obey
\begin{equation}\label{actual-top-lift-word}
 [\mathcal F_{l::w}\mathcal B_{l::w}]_{e,\mathbf0}
 =\left[\left(\sum_{i=0}^1 e_{i0}^{\rm new}
                           \mathcal C_{l,w,i}\right)
                       \mathcal F_w\mathcal B_w\right]_{e,\mathbf0},
\end{equation}
where $\mathcal C_{l,w}$ is the letter column just lifted.
This is ordinary two-by-two multiplication: the old frame
commutes with the site entries and is free at the new tensor
position, while the first column of the new site product is
$(A_r,B_r)$, or its original dual.  It introduces no extra
rational dressing.  Iterating
\eqref{actual-top-lift-compose}, the cost changes from
$(md,11mRd)$ to at most
\[
 ((m+1)d,\ 9Rd+11mRd+2Rd)
                       =((m+1)d,11(m+1)Rd).
\]
Starting with the identity proves \eqref{actual-top-lift-cost}
and the exact radial action.  This includes $r=0$ and the
empty word.  Extend the constructed coefficient fields by
zero off $\mathbf G^+$; no smoothness there is required of
the coefficients.  The original input cutoff makes every
resulting output smooth across that boundary.  Define
operators arbitrarily for non-cut words, where the assertion
imposes no condition, and impose the empty-word identity
globally.  Formula
\eqref{actual-top-lift-class} proves smoothness throughout
$\mathbf G^+$; its polynomial coefficient bounds and
$\mathrm{FibBound}$ prove preservation of the rapidly
decreasing class by the product rule.  This completes the
construction.\cite{top-lift-word-source}
\end{proof}

The theorem supplies an actual choice of the top-lift family
and its uniform differential complexity for every valid cut.
Together with the preceding terminal, Fourier and finite-decay
suppliers it removes this local analytic hypothesis.  Its
constant $11$ bounds differential order and numerator degree,
not the norms of the moving compact-type inverses.  The
arithmetic period maps, their integral lattices and the global
rational-return sums remain separate requirements.
